% Sakum Lang - Extension Reference (generated from EXTENSIONS.sakdoc)
\documentclass[11pt,a4paper]{article}
\usepackage[margin=2.2cm]{geometry}
\usepackage[T1]{fontenc}\usepackage[utf8]{inputenc}
\usepackage{longtable}\usepackage{array}\usepackage{xcolor}
\usepackage{hyperref}
\title{Sakum Lang \textemdash{} File-Type Extension Reference}
\author{Sakum Lang Project}
\date{Generated from \texttt{EXTENSIONS.sakdoc}}
\begin{document}
\maketitle

\section{Overview}
Every artifact in Sakum Lang carries a purpose-specific extension so the
compiler and tooling can recognise specialised data directly instead of
treating everything as plain text. The mapping is the \textbf{single source
of truth} loaded by \texttt{tools/sakum\_ext.py} and declared in
\texttt{sakum\_lang.sakproj}. Human-readable docs keep common formats
(\texttt{.md}, \texttt{.tex}, \texttt{.pdf}, \texttt{.html}).

\section{Extension Registry}
\begin{longtable}{|l|l|l|p{6.5cm}|}
\hline
\textbf{Extension} & \textbf{Category} & \textbf{Dispatch} & \textbf{Purpose} \\
\hline\endhead
\texttt{.sak} & source & build & Sakum source code \\
\hline
\texttt{.sakm} & source & build & Sakum module / package source \\
\hline
\texttt{.sakh} & source & build & Header / interface declarations \\
\hline
\texttt{.sakpkg} & package & manifest & Package manifest \\
\hline
\texttt{.sakproj} & package & manifest & Project configuration \\
\hline
\texttt{.saklock} & package & manifest & Dependency lock file \\
\hline
\texttt{.sakir} & ir & build & Intermediate Representation (SIR) \\
\hline
\texttt{.sakast} & ast & build & Abstract Syntax Tree \\
\hline
\texttt{.sakbc} & binary & link & Bytecode \\
\hline
\texttt{.sakobj} & binary & link & Object file \\
\hline
\texttt{.saklib} & binary & link & Static library \\
\hline
\texttt{.sakdll} & binary & link & Dynamic library (Windows) \\
\hline
\texttt{.sakso} & binary & link & Dynamic library (Linux) \\
\hline
\texttt{.sakdylib} & binary & link & Dynamic library (macOS) \\
\hline
\texttt{.sakexe} & binary & link & Platform-independent executable bundle \\
\hline
\texttt{.sakdoc} & doc & view & Language documentation \\
\hline
\texttt{.sakapi} & doc & view & API documentation \\
\hline
\texttt{.saktest} & test & test & Unit tests \\
\hline
\texttt{.sakbench} & test & test & Benchmarks \\
\hline
\texttt{.sakmath} & domain & view & Mathematical formulas / symbolic expressions \\
\hline
\texttt{.sakphys} & domain & view & Physics formulas \\
\hline
\texttt{.sakchem} & domain & view & Chemistry equations \\
\hline
\texttt{.sakbio} & domain & view & Biology / biotechnology models \\
\hline
\texttt{.sakquant} & domain & view & Quantum algorithms and circuits \\
\hline
\texttt{.sakml} & domain & view & Machine learning models / graphs \\
\hline
\texttt{.saktensor} & domain & view & Tensor expressions \\
\hline
\texttt{.sakproof} & domain & view & Formal proofs \\
\hline
\texttt{.sakgraph} & domain & view & Graph / network definitions \\
\hline
\texttt{.sakquery} & script & run & Query language scripts \\
\hline
\texttt{.sakschema} & data & validate & Data schemas \\
\hline
\texttt{.sakcfg} & config & validate & Configuration \\
\hline
\texttt{.sakcache} & cache & ignore & Compiler cache \\
\hline
\texttt{.sakidx} & index & ignore & Search / index database \\
\hline
\texttt{.sakdb} & data & validate & Embedded database \\
\hline
\texttt{.saklog} & log & view & Logs \\
\hline
\end{longtable}

\section{Reserved Core Set}
The following extensions are reserved for the language core and must not be
reused for other purposes:
\begin{center}
\texttt{.sak} \texttt{.sakir} \texttt{.sakast} \texttt{.sakmath} \texttt{.sakquant} \texttt{.sakdoc}
\end{center}

\section{Dispatch Semantics}
Behaviour is bound to the \textit{category}, not the file name. This is what
makes tooling honour the scheme uniformly:
\begin{itemize}
\item \texttt{build} -- compile through the pipeline (source $\to$ AST $\to$ IR $\to$ binary).
\item \texttt{link} -- feed to linker / loader.
\item \texttt{manifest} -- resolve dependencies / validate package.
\item \texttt{view} -- render to human-readable (doc / domain knowledge).
\item \texttt{test} -- run the test / benchmark harness.
\item \texttt{run} -- execute query scripts.
\item \texttt{validate} -- schema / config validation.
\item \texttt{ignore} -- derived data (cache / index); safe to drop.
\end{itemize}

\section{Generated Knowledge Files}
Domain knowledge is authored in dedicated extensions, e.g.:
\begin{itemize}
\item \texttt{Knowledge/Mathematics/\textasciitilde/Calculus.sakmath}
\item \texttt{Knowledge/Physics/\textasciitilde/mechanics.sakphys}
\item \texttt{Knowledge/Chemistry/\textasciitilde/organic.sakchem}
\item \texttt{Knowledge/Physics/Quantum Computing/gates.sakquant}
\end{itemize}
\end{document}